\section{Introduction}

% why is
One of the central promises of Unsupervised reinforcement learning (URL) is reusable-control: an agent first learns a task-agnostic representation of the environment, then uses it to solve multiple downstream tasks.
In robotics, for example, a robot may need to solve many downstream goals in the same workspace, each specified by a different reward function. Learning a new policy from scratch for every reward can be prohibitively expensive.
URL addresses this problem by first learning a representation, policy family, or value structure in a reward-free environment, and then using this object to infer policies for downstream rewards.

A wide range of URL algorithms instantiate this idea.
Goal-conditioned methods learn policies or values indexed by target states \citep{eysenbach-iclr21a,ma-neurips22a,park-iclr25a}.
Contrastive RL methods learn representations that separate future states from negatives \citep{eysenbach-neurips22a,wang-icml23a,park-iclr24a}.
Successor-feature based methods learn reward-free objects that make downstream policy inference tractable \citep{dayan-nc93a,barreto-neurips17a,touati-neurips21a,touati-iclr23a}.
Proto-successor measures, proto-value functions, skill-discovery methods, controllable representations, and world models make different assumptions about what should be represented during pretraining \citep{mahadevan-icml05a,mahadevan-jmlr07a,farebrother-iclr23a,eysenbach-iclr22a,agarwal-arxiv25a,efroni-iclr22a,lamb-arxiv22a,levine-arxiv24a,rudolph-rlc24a,hafner-iclr20a,hafner-arxiv23a}. Although these methods use different losses and inference rules, many of them reason about future state occupancy, state reachability, or action-dependent predictability in some form \citep{agarwal-rlj25a}.

An often underexplored area for such methods has been the domain of manipulation.
Cube manipulation is not merely navigation in a larger state space.
Moving a gripper to a cube, closing around it, lifting it, maintaining contact, and transporting it to a goal are different regimes of control.
The transition between regimes is mediated by contact.
Before grasp, gripper motion and object motion are weakly coupled; after grasp, they are tightly coupled; after a drop, the system returns to a different controllability regime.
A representation that captures coarse reachability can therefore be inadequate if it fails to preserve which phase the system is in and which actions maintain control of the object.

This paper investigates the failure mode of state-based URL policies in contact-rich cube manipulation.
We use manipulation as a diagnostic domain for understanding what existing representations preserve and what they miss.
Aggregate success is insufficient for this purpose.
A policy that never reaches the cube, a policy that reaches but cannot lift, and a policy that lifts but cannot maintain contact all receive the same final failure label, even though they fail for different algorithmic reasons.

% The hypothesis is phase-conditioned reachability.
\begin{hypothesis}[Contact bottleneck hypothesis]
Manipulation is hard for URL because the downstream task is not a single smooth reachability problem. It is a phase-conditioned control problem in which contact states act as bottlenecks. URL methods that compress future reachability without preserving contact phase may represent broad state reachability while failing to preserve the controllability needed for grasp maintenance and object transport.
\end{hypothesis}

% Contributions separate the formal diagnosis, the empirical protocol, and the empirical finding.
We present the diagnosis in three steps.
\begin{enumerate}
    \item We formulate manipulation as a phase-conditioned successor-measure problem and define a phase-consistency condition for state abstractions. This makes precise how phase-blind compression can become successor-inconsistent around contact bottlenecks.
    \item We introduce a state-only failure decomposition for \ogbench{} \task{}. Each rollout is assigned to the first failed physical precondition: approach, grasp, maintain grasp, or transport.
    \item We apply this diagnostic to \fb{} and \gciql{}. \fb{} reaches the cube in all analyzed rollouts and has fewer initial grasp failures than \gciql{}, but its aggregate deficit is concentrated after pick-up. A support diagnostic further indicates that the issue is not simply nearest-neighbor dataset coverage.
\end{enumerate}
